% Part: sets-functions-relations
% Chapter: relations
% Section: operations

\documentclass[../../../include/open-logic-section]{subfiles}

\begin{document}

\olfileid{sfr}{rel}{ops}
\olsection{સંબંધો પરની ક્રિયાઓ}

સંબંધોમાં ફેરફાર કરવો અથવા તેમને ભેગા કરવા ઘણી વાર ઉપયોગી
બને છે. \olref[sfr][rel][ord]{prop:stricttopartial}માં આપણે
સંબંધોના \emph{યોગગણ}નો વિચાર કર્યો હતો, જે બે સંબંધોને
જોડના ગણ તરીકે લઈએ ત્યારે તેમનો યોગગણ જ છે.
એ જ રીતે, \olref[sfr][rel][ord]{prop:partialtostrict}માં
આપણે સંબંધોના સાપેક્ષ તફાવતનો વિચાર કર્યો હતો.
અહીં સંબંધો પર કરી શકાય તેવી કેટલીક બીજી ક્રિયાઓ આપેલી છે.

\begin{defn}\ollabel{relationoperations}
ધારો કે $R$, $S$ સંબંધો છે, અને $A$ કોઈ પણ ગણ છે.

$R$નો \emph{વ્યસ્ત} એ
$R^{-1} = \Setabs{\tuple{y, x}}{\tuple{x, y} \in R}$ છે.

$R$ અને $S$નો \emph{સાપેક્ષ ગુણાકાર} એ
$(R \mid S) = \{\tuple{x, z} : \exists y(Rxy \land Syz)\}$ છે.

$R$નું $A$ પર \emph{મર્યાદન} એ
$\funrestrictionto{R}{A}= R \cap A^2$ છે.

$R$ને $A$ પર \emph{લાગુ પાડતાં મળતો ગણ} એ
$\funimage{R}{A} = \{y : (\exists x \in A)Rxy\}$ છે.
\end{defn}

\begin{ex}
ધારો કે $S \subseteq \Int^2$ એ $\Int$ પરનો અનુગામી સંબંધ છે,
એટલે કે $S = \Setabs{\tuple{x, y} \in \Int^2}{x + 1 = y}$,
જેથી $Sxy$ તો અને તો જ $x + 1 = y$.

$S^{-1}$ એ $\Int$ પરનો પુરોગામી સંબંધ છે,
એટલે કે $\Setabs{\tuple{x,y}\in\Int^2}{x -1 =y}$.

$S\mid S$ એ
$ \Setabs{\tuple{x,y}\in\Int^2}{x + 2 =y}$ છે.

$\funrestrictionto{S}{\Nat}$ એ $\Nat$ પરનો અનુગામી સંબંધ છે.

$\funimage{S}{\{1,2,3\}}$ એ $\{2, 3, 4\}$ છે.
\end{ex}

\begin{defn}[પરંપરિત સંવરણ]ધારો કે $R \subseteq A^2$ દ્વિઘટકી સંબંધ છે.

$R$નું \emph{પરંપરિત સંવરણ} એ
$R^+ = \bigcup_{0 < n \in \Nat} R^n$ છે,
જ્યાં આપણે પુનરાવર્તી રીતે $R^1 = R$ અને
$R^{n+1} = R^n \mid R$ વ્યાખ્યાયિત કરીએ છીએ.

$R$નું \emph{સ્વવાચક પરંપરિત સંવરણ} એ
$R^* = R^+ \cup \Id{A}$ છે.
\end{defn}

\begin{ex}
અનુગામી સંબંધ $S \subseteq \Int^2$ લો.
$S^2xy$ તો અને તો જ $x + 2 = y$,
$S^3xy$ તો અને તો જ $x + 3 = y$, વગેરે.
તેથી $S^+xy$ તો અને તો જ કોઈક $n \geq 1$ માટે
$x + n = y$. બીજા શબ્દોમાં, $S^+xy$ તો અને તો જ $x < y$,
અને $S^*xy$ તો અને તો જ $x \le y$.
\end{ex}

\begin{prob}
$R$નું પરંપરિત સંવરણ ખરેખર પરંપરિત છે તે દર્શાવો.
\end{prob}

\end{document}
